\documentclass[10.5pt]{beamer}

%\linespread{1.2}
%    Include other referenced packages here.
\usepackage{pgfplots}
\usepackage{mathtools}
\usepackage{amssymb,amsmath,amsthm}
%\usepackage[inline]{enumitem}
\usepackage{relsize}
\usepackage{mathtools}
\usepackage{multicol}
\usepackage{booktabs}
\usepackage{calc}
\usepackage{esdiff}
\usepackage{systeme}
\usepackage{bm}
\usepackage[most]{tcolorbox}

%%%%%%%%%%%%%%%%%%%%%%%%%%%%%%%%%%%%%%%%%%%%%%%%%%%%%%%%%%%%%%%%%%%%%%%%
% definitions
% theorem, proposition, corollary, lemma
% \newtheorem{theorem}{Theorem}[section]
% \newtheorem{corollary}{Corollary}[theorem]
% \newtheorem{lemma}[theorem]{Lemma}
% \newtheorem{proposition}{Proposition}[section]
% \newtheorem{definition}{Definition}[section]


\newcommand{\mydef}[2]{
    \begin{tcolorbox}[colback=blue!5, colframe=blue!75!black, title=#1]
        #2
    \end{tcolorbox}
}


% numerical sets
\newcommand{\real}{{\mathbb R}}
\newcommand{\realm}{{{\mathbb R}^m}}
\newcommand{\rational}{{\mathbb Q}}
\newcommand{\integer}{{\mathbb Z}}
\newcommand{\naturals}{{\mathbb N}}
\newcommand{\complex}{{\mathbb C}}
\newcommand{\F}{{\mathbb F}}
\newcommand{\R}{{\mathbb R}}
\newcommand{\C}{{\mathbb C}}
\newcommand{\Rn}{{\mathbb R}^n}
\newcommand{\Rm}{{\mathbb R}^m}
\newcommand{\cn}{{\mathbb C}^n}
\newcommand{\rnn}{{{\mathbb R}^{n\times n}}}
\newcommand{\rnm}{{{\mathbb R}^{n\times m}}}
\newcommand{\rmn}{{{\mathbb R}^{m\times n}}}
\newcommand{\cnn}{{{\mathbb C}^{n\times n}}}

% abs and norm, ceil and floor, inner product
\DeclarePairedDelimiter\abs{\lvert}{\rvert}%
\DeclarePairedDelimiter\norm{\lVert}{\rVert}%
\DeclarePairedDelimiter\ceil{\lceil}{\rceil}
\DeclarePairedDelimiter\floor{\lfloor}{\rfloor}
\DeclarePairedDelimiter\inner{\langle}{\rangle}
%\providecommand{\norm}[1]{\lVert#1\rVert}
% Swap the definition of \abs* and \norm*, so that \abs
% and \norm resizes the size of the brackets, and the 
% starred version does not.
\makeatletter
\let\oldabs\abs
\def\abs{\@ifstar{\oldabs}{\oldabs*}}
%
\let\oldnorm\norm
\def\norm{\@ifstar{\oldnorm}{\oldnorm*}}
\makeatother

% common functions
\newcommand{\rtr}{\real\to\real}
\newcommand{\rttr}{\real^2\to\real}

% common matrices
\newcommand{\MmnF}{M_{mn}(\F)}
\newcommand{\MnF}{M_{n}(\F)}

% diagonal matrix
\DeclareMathOperator{\diag}{diag}

%domain 
\DeclareMathOperator{\dom}{dom}

% image of a linear operator
\DeclareMathOperator{\Ima}{Im}

% range of a linear operator
\DeclareMathOperator{\range}{range}

% Rank of a linear operator
\DeclareMathOperator{\rank}{rank}

% projection of a linear operator
\DeclareMathOperator{\proj}{proj}

% Null and Column space of a matrix
\DeclareMathOperator{\rspace}{\mathcal{R}}
\DeclareMathOperator{\nspace}{\mathcal{N}}
\DeclareMathOperator{\cspace}{\mathcal{C}}

% Trace of a linear operator
\DeclareMathOperator{\tr}{tr}

% Rank of a linear operator
\DeclareMathOperator{\vspan}{span}

% Sign of a function
\DeclareMathOperator{\sign}{sgn}

% differential d
\newcommand{\dif}{\mathop{}\!\mathrm{d}}
\newcommand{\od}[2]{\frac{\dif #1}{\dif #2}}
\newcommand{\dod}[2]{\displaystyle\frac{\dif #1}{\dif #2}}

% complement
\newcommand{\scomp}[1]{{#1^\complement}}

% interior & closure
\newcommand{\interior}[1]{{\kern0pt#1}^{\mathrm{o}}}
\newcommand{\closure}[1]{{\overline{#1}}}

% evaluated integral 
\newcommand\eval[3]{\left[#1\right]_{#2}^{#3}}

% matrices p q
\newcommand{\rmat}[2]{{{\mathbb R}^{{#1}\times {#2}}}}

% vectors
\newcommand{\vect}[1]{\mathbf{#1}}
\newcommand{\va}{\mathbf{a}}
\newcommand{\vb}{\mathbf{b}}
\newcommand{\vc}{\mathbf{c}}
\newcommand{\vd}{\mathbf{d}}
\newcommand{\ve}{\mathbf{e}}
\newcommand{\vf}{\mathbf{f}}
\newcommand{\vp}{\mathbf{p}}
\newcommand{\vq}{\mathbf{q}}
\newcommand{\vr}{\mathbf{r}}
\newcommand{\vs}{\mathbf{s}}
\newcommand{\vx}{\mathbf{x}}
\newcommand{\vy}{\mathbf{y}}
\newcommand{\vz}{\mathbf{z}}
\newcommand{\vu}{\mathbf{u}}
\newcommand{\vw}{\mathbf{w}}
\newcommand{\vv}{\mathbf{v}}
\newcommand{\vT}{\mathbf{T}}
\newcommand{\vX}{\mathbf{X}}
\newcommand{\vY}{\mathbf{Y}}

\newcommand{\vzero}{\mathbf{0}}
\newcommand{\mzero}{\mathbf{O}}

\newcommand{\vtn}[4]{
\vect{#1}=\left[
\begin{array}{c}
#2 \\
#3 \\
#4
\end{array}
\right]}

\newcommand{\vtu}[3]{
\left[
\begin{array}{c}
#1 \\
#2 \\
#3
\end{array}
\right]}

\newcommand{\vtd}[4]{
\left[
\begin{array}{c}
#1 \\
#2 \\
#3 \\
#4
\end{array}
\right]}

\newcommand{\vdn}[3]{
\vect{#1}=\left[
\begin{array}{c}
#2 \\
#3 
\end{array}
\right]}

\newcommand{\vdu}[2]{
\left[
\begin{array}{c}
#1 \\
#2 
\end{array}
\right]}

\newcommand{\stlb}[3]{
\left\{\begin{matrix}
#1 \\ #2 \\ #3
\end{matrix}\right.}

\newcommand{\stlt}[3]{
\left\{\begin{matrix}
#1 \\ #2 \\ #3
\end{matrix}\right.}

% horbars in matrices
\newcommand{\brows}[1]{%
  \begin{bmatrix}
  \begin{array}{@{\protect\rotvert\;}c@{\;\protect\rotvert}}
  #1
  \end{array}
  \end{bmatrix}
}
\newcommand{\rotvert}{\rotatebox[origin=c]{90}{$\vert$}}
\newcommand{\rowsvdots}{\multicolumn{1}{@{}c@{}}{\vdots}}


\usepackage{bm} 
\usepackage{amsmath}
\usepackage{xcolor}
\usepackage{comment}
\usepackage{graphicx}
\usepackage{listings}
\setbeamertemplate{caption}[numbered]
\usetheme{CambridgeUS}
\usecolortheme{dolphin}
\usepackage{color}
\usepackage{subcaption}
\usepackage[most]{tcolorbox}
\usepackage{tikz}
\usepackage{pifont}

\usepackage[round]{natbib} 
\bibliographystyle{plainnat} 

\setbeamercovered{transparent}

\definecolor{dkgreen}{rgb}{0,0.6,0}
\definecolor{gray}{rgb}{0.5,0.5,0.5}
\definecolor{mauve}{rgb}{0.58,0,0.82}

\lstset{frame=tb,
  language=Python,
  aboveskip=3mm,
  belowskip=3mm,
  showstringspaces=false,
  columns=flexible,
  basicstyle={\small\ttfamily},
  numbers=none,
  numberstyle=\tiny\color{gray},
  keywordstyle=\color{blue},
  commentstyle=\color{dkgreen},
  stringstyle=\color{mauve},
  breaklines=true,
  breakatwhitespace=true,
  tabsize=3
}

\newcommand{\methoditem}[2]{
\ifnum#1=\value{methodindex}
    {\color{blue}\textbf{$\blacktriangleright$ #2}}
\else
    {\color{gray}#2}
\fi
}


\newcounter{methodindex}

\setbeamersize{text margin left=8mm,text margin right=8mm} 
\beamertemplatenavigationsymbolsempty

\hypersetup{
    colorlinks = true,
    linkcolor = blue!60!black,
    citecolor = red
}

\setbeamertemplate{footline}{
  \hfill%
  \usebeamercolor[fg]{page number in head/foot}%
  \usebeamerfont{page number in head/foot}%
  \setbeamertemplate{page number in head/foot}[framenumber]%
  \usebeamertemplate*{page number in head/foot}\kern1em\vskip2pt%
}

\AtBeginSection[]
{
  \begin{frame}
    \frametitle{Table of Contents}
    \tableofcontents[currentsection]
  \end{frame}
}

\title{Network Independence Testing via Copulas}
\author{\vspace{-0.5cm} Lorenzo Costa}
\institute{}
\date{March 2026}

\begin{document}

\begin{frame}
\titlepage
\end{frame}

% \begin{frame}{Contents}
% 	\tableofcontents
% \end{frame}

\begin{frame}{Problem Definition}
\textbf{Goal}: test the independence between two layers of a multiplex network. For adjacency matrices $A\in \mathbb{R}^{n\times n}, B\in \mathbb{R}^{n\times n}$ we want to test:
\[
H_0: A\perp B \text{ vs } H_1: A\not\perp B
\]

\textbf{Setup}: Assume there exists latent variables $z_{i}, x_{i}\in \mathbb{R}^d$ such that entries of the adjacency matrices are \textbf{conditionally independent} given these latent variables:
\[
p(A_{ij}, B_{ij}) = p_A(A_{ij}|z_{i}, z_{j}) p_B(B_{ij}|x_{i}, x_{j}) \;p_{z, x}(z_{i}, z_{j}, x_{i}, x_{j})
\]
\end{frame}

% \begin{frame}{Why do we care ?}
    
% \end{frame}

\begin{frame}{Why do we need this?}
For example in economics understanding systemic financial risk requires analyzing how institutions interact across different markets. Modelling these as multiple networks on the same nodes relies on a key assumption:
\begin{itemize}
    \item if they are \textbf{independent} we can just sum the individual risks
    \item if they are \textbf{dependent}  summing individual risk leads to significant under-estimation of risk e.g. \citep{poledna_multi-layer_2015} 
\end{itemize}

\begin{figure}
    \centering
    \includegraphics[width=0.5\linewidth]{figures/risk.png}
    \caption{Risk for different markets - Mexican Banking system}
    \label{fig:placeholder}
\end{figure}

\end{frame}

\section{From the literature}

\begin{frame}{From the literature}
From the literature two main threads of research emerge:
\begin{enumerate}
\setlength{\itemsep}{1.2em}

    \item tests based on \textbf{latent positions} such as based on model assumption \citep{fosdick_testing_2015}, distance correlation \citep{lee_network_2019} or RKHS \citep{wen_rkhs-based_2023}.
    \begin{itemize}
        \item \textbf{cons}: requires estimation of latent positions and choice of embedding dimension
        \item \textbf{pros}: reduces noise, possibly higher power
    \end{itemize}
    \pause
    \item tests based on the \textbf{adjacency matrix} itself such as the Quadratic Assignment Procedure  \citep{mantel_detection_1967, shi_asymptotic_2025}
    \begin{itemize}
        \item \textbf{cons}: more sensitive to noise, may have lower power.
        \item \textbf{pros}: no embedding needed, easier to use 
    \end{itemize}
    
\end{enumerate}

\end{frame}

\begin{frame}{From the literature}
All the methods presented later have a similar structure:
\begin{block}{Testing Procedure}
\begin{itemize}
    \item \textbf{Step 1: embedding} (optional). Estimate latent positions $\hat{X}, \hat{Z}$.
    \item \textbf{Step 2: test statistic}. Compute dependence some measure of dependence.
    \item \textbf{Step 3: null distribution}. Approximate the null distribution by permuting the nodes and obtain a p-value.
\end{itemize}
\end{block}
By choosing step 1 and 2 we obtain different tests having specific strengths and weaknesses.
\end{frame}

\subsection{Distance Correlation}

\begin{frame}{From the literature - Distance Covariance}
Define $(Z', X')$ iid copies of random vectors $(Z, X)$ and:
\[
m_1(X, x) =\mathbb{E}_{X}||X-x|| \quad m_2(X, X') = \mathbb{E}_{X, X'}||X-X'||
\]

\mydef{Distance Covariance \citep{szekely_measuring_2007}}{
Distance covariance is defined as:
\[
dcov(Z, X) = \mathbb{E}_{Z, Z', X, X'}(d(Z, Z')d(X, X')) 
\]
for $d(x, x') = ||x-x'|| - m_1(X, x) - m_1(X', x') + m_2(X, X')$. 
}
Distance correlation is then defined by dividing $dcov$ by distance variances.

$\Rightarrow$ it can be shown that $dcov$ is zero if and only if $Z, X$ are independent 
\end{frame}

\begin{frame}{From the literature - Distance Correlation}
We can define a sample version as:
\[
dcov(Z, X) = \frac{1}{n^{2}}\sum_{i,j}\tilde{C}_{ij}\tilde{D}_{ij}
\]
for $\tilde{D}$ the doubly centred matrix of distances $D_{ij}=||Z_i - Z_j||$ (and $\tilde{C}$ in a similar way for $X$)
\vspace{.5em}

\citet{shen_distance_2020} further extend this to a "local optimal" version, called \textbf{Multiscale Graph Correlation} (MGC), which has higher finite-sample power. 

\end{frame}

\begin{frame}{Multiscale Graph Correlation}
Define $N_{i, k}$ as the set containing the $k$ closest nodes to $i$ and the truncated matrix $D_{ij}^{(k)}=D_{ij}1(j \in N_{i, k})$ (and $\tilde{D}^{(k)}$ its centred version).
\mydef{Multiscale Graph Correlation}{
Define local distance covariance as:
\vspace{-.5em}
\[
dcov^{k,m}(Z, X) = \frac{1}{n^{2}}\sum_{i,j}\tilde{C}^{(k)}_{ij}\tilde{D}^{(m)}_{ij}
\] 
\vspace{-1.5em}

and MGC as the maximum local correlation:
\vspace{-.5em}
\[
MGC(Z, X) = \max_{k, m} dcorr^{k, m}(Z, X)
\]
} 
\end{frame}


\begin{frame}{From the literature - Distance Correlation}
\citet{lee_network_2019} present a method to test independence between two networks that is:
\begin{itemize}
    \item \textbf{valid} if the two networks are exchangeable 
    \item consistent under \textbf{any} form of dependency if $A, B$ are constrained to $\Theta$ where the embedding map is injective
\end{itemize}

\begin{block}{Testing Procedure }
\begin{itemize}
    \item \textbf{Step 1}: estimate latent positions $\hat{X}, \hat{Z}$ using \textit{Diffusion Maps} \citep{coifman_geometric_2005}
    \item \textbf{Step 2}: compute dependence measures $MGC(\hat{X}, \hat{Z})$ or $dCorr(\hat{X}, \hat{Z})$.
    \item \textbf{Step 3}: approximate the null distribution by permuting the rows of $\hat{X}$
\end{itemize}
\end{block}

% Limitations of this approach:
% \begin{itemize}
%     \item \textbf{practical}: step 1 requires estimation of the latent positions
%     \item \textbf{theoretical}: distance correlation works under the assumption that the latent positions have finite expectation
% \end{itemize}
\end{frame}

% \begin{frame}{From the literature - RKHS}
% Assume networks are generated as:
% \[
% A_{ij}\sim Bernoulli(K_{z_i,z_j}) \quad B_{ij}\sim Bernoulli(L_{x_i, x_j})
% \]
% for $K_{z_i,z_j}=K(z_i, z_j), L_{x_i,x_j}=L(x_i, x_j)$ and $K(\cdot, \cdot), L(\cdot, \cdot)$ kernel functions.
% \vspace{.5em}

% By Mercer's theorem we can write $K$ in its feature map representation:
% \[
% K(x, y) = \langle \phi(x), \phi(y) \rangle_{H} \quad \phi(x) = \{ \sqrt{\lambda_j}\psi_j(x), j=1, 2, \dots\}
% \]

% Denote as $\phi_d(x)$ the truncation of the feature map at $j=d$. Using a result in \citet{tang_universally_2013} we have that ASE 
% can be used to estimate the $\phi_d(x)$ 
% %with error decreasing as $O(\sqrt{\log n/ n})$ 
% such that the kernel matrix can be estimated as $\hat{K}=\hat{\phi}_d \hat{\phi}_d^T$.
% \end{frame}

% \begin{frame}{From the literature - RKHS}
% For $(X', Z'), (X'', Z'')$ independent copies of $(X, Z)$ \citet{wen_rkhs-based_2023} define \textit{Graph Covariance} as:
% \[
% gCov(A, B) = \mathbb{E}K_{XX'}L_{ZZ'} + \mathbb{E}K_{XX'}\mathbb{E}L_{YY'} - 2\mathbb{E}K_{XX'}L_{ZZ''}
% \]
% which is zero iff $G_1, G_2$ are independent.
% \vspace{.5em}

% A sample version can be found from the $U$-centred estimated matrices $\tilde{K}, \tilde{L}$ as:
% \[
% gCov_n(A, B) = \frac{1}{n(n-3)}\sum_{i\neq j} \tilde{k}_{ij} \tilde{l}_{ij}
% \]
% Using this measure of dependence the authors define a permutation test in the same way as for distance correlation.

% \end{frame}

\subsection{Quadratic Assignment Procedure}

\begin{frame}{From the literature - QAP}
QAP is a general framework to test the relationships between two networks using a double permutation strategy: 
\begin{block}{Quadratic Assignment Procedure}
\begin{itemize}
\setlength{\itemsep}{.5em}
    \item \textbf{Step 1}: compute a test statistic $W(A, B)$ 
    \item \textbf{Step 2}: permute rows and columns of $A$ according to permutation $\pi$ and recomputing $W^\pi(A_\pi, B)$ to get the p-value
\end{itemize}
\end{block}
\pause

A common choice of test statistic is the Pearson correlation:
\[
\hat{\rho}(A, B) = \frac{\sum_{i\neq j}(A_{ij}-\bar{A})(B_{ij}-\bar{B})}{\sqrt{\sum_{i\neq i}(A_{ij}-\bar{A})^{2}\sum_{i\neq j}(B_{ij}-\bar{B})^{2}}}
\]
Which would be the equivalent of computing $dcorr$ using L2 distance on the observed adjacency matrix.
\end{frame}

\begin{frame}{From the literature - QAP}
Assume now that there exists two symmetric functions $\alpha(\cdot, \cdot), \beta(\cdot, \cdot)$ such that entries of the adjacency matrices can be written as:
\[
A_{ij} = \alpha(Z_i, Z_j)\quad B_{ij}= \beta(X_i, X_j)
\]
For a general kernel $\phi$ define $\phi_0=\mathbb{E}\phi(X, X')$ and:
\[
\underbrace{\eta_{\phi, 1} = \mathbb{E}_X\left(\mathbb{E}_{X'}\phi(X, X')-\phi_0 \right)^2}_{\text{first order variance}} \quad \underbrace{\eta_{\phi, 2} = \mathbb{E}_{X, X'}\left(\phi(X, X')-\phi_0 \right)^2}_{\text{second order variance}}
\]
\citet{shi_asymptotic_2025} show that, under $H_0: Z\perp X$, the sampling distribution of $\hat{\rho}$ is:
\[
\sqrt{n}\hat{\rho}\to \mathcal{N}\left(0, v\right) \quad v = \frac{\eta_{1, \alpha}\eta_{1, \beta}}{\eta_{2, \alpha}\eta_{2, \beta}}
\]

\end{frame}

\begin{frame}{From the literature - QAP}
Moreover assuming that there exists two functions $\kappa_\alpha(\cdot), \kappa_\beta(\cdot),$ such that:
\[
|\alpha(Z, Z')| \le \kappa_\alpha(Z) + \kappa_\alpha(Z') \quad |\beta(X, X')| \le \kappa_\beta(X) + \kappa_\beta(X') 
\]
and that $\kappa_\alpha(\cdot), \kappa_\beta(\cdot),$ have finite fourth moment, we have:
\[
\lim_{n\to \infty}\sup_{t \in \mathbb{R}} |\mathcal{L}(t; \sqrt{n}\hat{\rho}^\pi) - \mathcal{L}(t; \mathcal{N}(0, v)|=0 \quad a.s.
\]
where $\mathcal{L}(t; \sqrt{n}\hat{\rho}^\pi)=\frac{1}{n!}\sum_{\pi\in \mathbb{S}_n} 1(\sqrt{n}\hat{\rho}^\pi \le t)$ denotes the permutation cdf of $\hat{\rho}$ computed using all $n!$ permutations
\vspace{.5em}
% \pause
% \begin{itemize}
%     \item[{\ding{226}}] \textbf{pros}: very easy to use
%     \item[{\ding{226}}] \textbf{cons}: can detect only monotonic dependence
% \end{itemize}

\end{frame}

\section{Copula perspective}
\begin{frame}{What is a Copula?}
Take a random vector $(X_1, X_2, \dots, X_d)$ with continuous marginals. By applying the probability integral transform to each component, the random vector:
\[
(U_1, U_2, \dots, U_d) = (F_1(X_1), F_2(X_2), \dots, F_d(X_d))
\]
has marginals that are uniform on $[0, 1]$.
\vspace{.5em}

The copula of $(X_1, \dots, X_d)$ is defined as the joint cdf of $(U_1, \dots, U_d)$:
\[
\begin{split}
    C(u_1, u_2, \dots, u_d) &= \mathbb{P}[U_1 \leq u_1, U_2 \leq u_2, \dots, U_d \leq u_d] \\
    &= \mathbb{P}[X_1 \leq F_1^{-1}(u_1), X_2 \leq F_2^{-1}(u_2), \dots, F_d^{-1}(u_d),]
\end{split}
\]
$\Rightarrow$ this contains all information on the dependence structure between the components of $(X_1, X_2, \dots, X_d)$.
\end{frame}

\begin{frame}{Framework set-up}
Assuming the pairs $(x_i, z_i)$ are iid we can write: we can write the likelihood:
\vspace{-.5em}
\[
p(A_{ij}, B_{ij}) = p_A(A_{ij}|z_{i}, z_{j}) p_B(B_{ij}|x_{i}, x_{j}) \;p_{Z, X}(z_{j},x_{j})p_{Z, X}(z_{i}, x_{i})
\]
\vspace{-1.25em}

By Sklar’s theorem \citep{sklar_fonctions_1959} any multivariate cdf can be decomposed into its \textcolor{green!50!black}{marginals} and a \textcolor{blue}{copula}\footnote{here we're assuming the density exists, for discrete distributions copula still exists but $p_{x,z}$ is more messy}:
\[
p_{z, x}(z_{i}, x_{i}) = {\color{green!50!black}f_z(z_{i, 1}) \dots f_z(z_{i, d}) f_x(x_{i, 1}) \cdot f_x(x_{i, d})} \cdot {\color{blue}c_{z, x}}
\]
where ${\color{blue}c_{z, x}}= c(F_{z_{1}}(z_{i, 1}), \dots, F_z(z_{i,d}), F_x(x_{i,1}), \dots,,F_{x_{1}}(x_{i, d}) )$ denotes the copula density (i.e. derivative of copula)
\vspace{.5em}
\pause

If we further assume that the dimensions of the latent space are independent we get:
\vspace{-.5em}
\[
p_{z, x}(z_{i}, x_{i}) = \prod_{k=1}^d {\color{green!50!black}f_z(z_{i, k})f_x(x_{i, k})}{\color{blue}c(F_z(z_{i, k}), F_x(x_{i, k}))}
\]
% \vspace{-1.25em}

% which means we are in practice using a series of bivariate not high-dimensional copulas.
\end{frame}


\subsection{Copula constructions}
\begin{frame}{Copula constructions}
Two popular ways to construct copulas are:
\begin{itemize}
    \item \textbf{Elliptical copulas} which are built from well-known multivariate distributions such as the Gaussian distribution as:
    \[
    C_{gauss}(\mathbf{u}) = \Phi_{R}(\Phi_{1}^{-1}(F_{1}(u_1)),\dots, \Phi_{1}^{-1}(F_{d}(u_d)))
    \]
    for $\Phi_R$ the cdf of a multivariate Gaussian with correlation matrix $R$ and $\Phi_1$ the cdf of $\mathcal{N}(0, 1)$.
    \item \textbf{Archimedean copulas} which are build starting from a function $\phi(\cdot)$ (called a generator) as:
    \[
    C_{arch}(u, v) = \phi^{[-1]}(\phi(u)+\phi(v))
    \]
    For $\phi(t): [0,1]\to[0, \infty)$ a strictly decreasing, continuous, and convex function with $\phi(1)=0$.
\end{itemize}


\vspace{.5em}

\end{frame}

\begin{frame}{Copula constructions}
Different choices of $\phi(\cdot)$ allow us to model many kind of dependencies such as:
\begin{itemize}
    % \item \textbf{Frank} copula: symmetric with no tail dependency. Generator is $\phi(u) = - \ln \left( \frac{e^{-u\theta}-1}{e^{-\theta}-1} \right)$ for $\theta\neq 0$
    \item \textbf{Clayton} copula: lower tail dependency. Generator is $\phi(u) = \frac{1}{\theta} (u^{-\theta}-1)$ for $\theta>0$
    \item \textbf{Gumbel} copula: upper tail dependency. Generator is $(-\ln (u))^\theta$ for $\theta\ge 1$.
\end{itemize}
\vspace{.5em}

Note: also that convex linear combination of copulas yield valid copulas so we can get highly non-linear relations by defining a \textbf{mixture of copulas} (e.g. mixture of Gaussian):
\[
C(u, v) = w_{1} C^{gauss}_{1}(u, v) + w_2C^{gauss}_{2}(u, v)
\]
\end{frame}

\begin{frame}<1-3>[t]{Copula constructions} % <1-3> forces the slides to exist
Different combinations of copula and marginals give us different correlation structures:

\begin{tikzpicture}[remember picture, overlay]
    % 1. Image on ALL slides (1, 2, and 3)
    % Using yshift to nudge it down if the title blocks it
    \node<1-> at ([yshift=-1cm]current page.center) {
        \includegraphics[width=1\textwidth]{figures/visualise_latent.png}
    };

    % 2. Text Box ONLY on slide 2
    \node<2> [
        fill=white,
        draw=black,
        line width=1pt, 
        inner sep=15pt,
        rounded corners,
        fill opacity=0.9, % Makes it slightly see-through so it looks "on top"
        text opacity=1     % Keeps text solid
    ] at (current page.center) {
        \begin{minipage}{0.75\textwidth}
            Note: for Elliptical and Archimedean copula there is a map between the copula parameter and Kendall's $\tau$:
            \[
            \tau_{\text{arch}} = 1+4 \int_0^1 \frac{\phi(t)}{\phi'(t)}dt 
            \quad \tau_{\text{ellip}} = \frac{2}{\pi}\arcsin(\rho)
            \]
        \end{minipage}
    };
\end{tikzpicture}
\end{frame}

\subsection{RV coefficient}
\begin{frame}{RV coefficient}
Under a Gaussian copula model testing independence reduces to testing for zero correlation. This motivates the use of a "matrix version" of Pearson correlation to build a test

\mydef{RV coefficient \citep{escoufier_traitement_1973, robert_unifying_1976}}{
Suppose $Z, X\in \mathbb{R}^{n\times d}$ are matrices with centred columns. For $\Sigma_{zx} =\mathbb{E}(ZX^T), \Sigma_{zz}=\mathbb{E}(ZZ^T)$, RV coefficient is defined as:
\[
RV(Z, X) = \frac{Tr(\Sigma_{zx}\Sigma_{xz})}{\sqrt{Tr(\Sigma_{xx})Tr(\Sigma_{zz})}}
\]
}
\end{frame}

\begin{frame}{RV Coefficient}
\vspace{0.5em}
\citet{mordant_measuring_2021} show that the maximum value the denominator can take is $\sum_{i=1}^d \lambda_{i, x}\lambda_{i, z}$ which motivates the definition of the \textbf{adjusted RV coefficient}:
\[
 RV(Z, X) = \frac{Tr(\Sigma_{ZX}\Sigma_{XZ})}{\sum_{i=1}^d \lambda_{i, z}\lambda_{i, x}}
\]
for $\lambda_{i, z}\lambda_{i, x}$ eigenvalues of $\Sigma_{ZZ}$ and $\Sigma_{XX}$.
\vspace{.5em}

We can then find a sample version, for $\hat{Z}, \hat{X}$ matrix of estimated latent positions, as:
\[
RV(\hat{Z}, \hat{X}) = \frac{Tr(\hat{Z}\hat{X}^T\hat{X}\hat{Z}^T)}{\sum_{i=1}^d \lambda_{i, z}\lambda_{i, x}}
\]
for $\lambda_{i, z}\lambda_{i, x}$ eigenvalues of $ZZ^T$ and $XX^T$.
\end{frame}

\begin{frame}{RV Coefficient}
\begin{block}{Testing procedure}
\begin{enumerate}
    \item \textbf{Step 1}: estimate the latent positions $\hat{X}, \hat{Z}$ 
     \item \textbf{Step 2}: compute adjusted RV coefficient
     \item \textbf{Step 3}: approximate the null distribution by permuting the rows of $\hat{X}$
\end{enumerate} 
\end{block}
\end{frame}

\begin{frame}{RV coefficient - when does it work?}
Under a Gaussian copula model testing independence reduces to testing for zero correlation. 
\vspace{.5em}

Moreover, using Hoeffings covariance decomposition:
\[
Cov(Z, X) = \int_{-\infty}^{\infty}\int_{-\infty}^{\infty} C(F_Z(z), F_X(x)) - F_Z(z)F_X(x) dz dx
\]
we can show that the equivalence between testing for zero correlation and independence, holds for any copula for which:
\begin{itemize}
    \item strict \textbf{Positive or Negative Quadrant dependence} holds (i.e. $C(u, v) - u \cdot v$ strictly greater or smaller than zero) 
    \item variance of the marginals exists
\end{itemize}
\vspace{.5em}

Examples of this include one parameter Archimedean Copulas and elliptical copulas
%TODO make this more general using multivariate Hoeffding lemma Block & Fang (1988) 
\end{frame}

\section{Possible extensions}
\begin{frame}{Desiderata}
We have that:
\begin{itemize}
    \item QAP is easy to use but cannot detect all forms of dependence
    \item distance correlation is able to detect all types of dependence but requires embedding the network
\end{itemize}
\vspace{.5em}

We are then missing something that:
\begin{itemize}
    \item can be applied "out-of-the box"
    \item requires minimal effort to choose hyper-parameters
    \item can detect all forms of correlation
\end{itemize}
\end{frame}

\subsection{Cramer Von Mises}
\begin{frame}{Cramer Von Mises}
\mydef{Cramer Von Mises \citep{jing_blum-kiefer-rosenblatt_1996}}{
For $F_X(\cdot), F_Z(\cdot)$ cdf of $X, Z$ and $F_{Z, X}(\cdot, \cdot)$ define:
\vspace{-.5em}
\[
T(Z, X)=\int (F_{Z, X}(z, x) - F_Z(z)F_X(x))^2 dzdx
\]
\vspace{-1.5em}

This is zero if and only if $Z, X$ are independent.
}

\pause 

Note that we can't apply this to latent positions because of rotational invariance.
\vspace{.5em}

However we can define a joint distribution on the \textbf{inner products} $P_{ij}^A=z_i^T z_j=\mathbb{E}(A_{ij})$ and apply CvM on these:
\[
T_n(A, B)=\int \Big(F_{A, B}(P_{ij}^A, P_{ij}^B) - F_A(P_{ij}^A)F_B(P_{ij}^B)\Big)^2 \, dF_AdF_B
\]
% $\Rightarrow$ the issue with applying this approach to latent variables is the orthogonal invariance if $d>2$.
\end{frame}

\begin{frame}{Cramer Von Mises}
For $F_n(\cdot)$ the empirical cdf of $P_{ij}$ we can get a plug-in estimator as:
\[
T_n(A, B) =\frac{1}{n(n-1)} \sum_{i\neq j} \Big(F_{n}(P_{ij}^A, P_{ij}^B) - F_{n}(P_{ij}^A)F_{n}(P_{ij}^B)\Big)^2
\]
Now defining as $U^A_{ij},U^B_{ij}$ the ranks of $P^A_{ij}, P^B_{ij}$ (i.e. ordering) we can get a copula version that strips the effect of the marginals as:
\[
T_n(A, B) =\frac{1}{n(n-1)} \sum_{i\neq j} \Big(\hat{C}_n(U^A_{ij}, U^B_{ij}) - U^A_{ij}U^B_{ij}\Big)^2
\]
for $\hat{C}(u^A, u^B) = \frac{1}{n(n-1)}\sum_{i\neq j}1(U^A_{ij}\le u^A, U^B_{ij}\le u^B)$ the empirical copula.

\end{frame}

% \begin{frame}{Cramer Von Mises}
% We can define a joint distribution on the \textbf{inner products} $z_i^T z_j=\mathbb{E}(A_{ij})$ and apply CvM on these:
% \[
% T_n(A, B)=\int \Big(F_{Z, X}(z_i^T z_j, x_i^T x_j) - F_Z(z_i^T z_j)F_X(x_i^T x_j)\Big)^2 \, dF_AdF_B
% \]
% For $F_n(\cdot)$ the empirical cdf of $P_{ij}$ we can get a plug-in estimator as:
% \[
% T_n(A, B) =\frac{1}{n(n-1)} \sum_{i\neq j} \Big(F_{n}(P_{ij}^A, P_{ij}^B) - F_{n}(P_{ij}^A)F_{n}(P_{ij}^B)\Big)^2
% \]
% Now defining as $U^A_{ij},U^B_{ij}$ the ranks of $P^A_{ij}, P^B_{ij}$ (i.e. ordering) we can get a copula version that strips the effect of the marginals as:
% \[
% T_n(A, B) =\frac{1}{n(n-1)} \sum_{i\neq j} \Big(\hat{C}_n(U^A_{ij}, U^B_{ij}) - U^A_{ij}U^B_{ij}\Big)^2
% \]
% for $\hat{C}(u^A, u^B) = \frac{1}{n(n-1)}\sum_{i\neq j}1(U^A_{ij}\le u^A, U^B_{ij}\le u^B)$ the empirical copula.

% \end{frame}

\begin{frame}{Ranking estimation - Weighted network}

Note that we don't need to estimate $P_{ij}$ but only their \textbf{ranks}. 
Take a weighted network:
\[
A_{ij} \overset{iid}{\sim}\mathcal{N}(z_i^T z_j, \sigma^2)
\]
or a RDPG model:
\[
A_{ij}\sim Bernoulli(z_i^T z_j) \quad z_i
\sim F \; s.t. \; z_i^T z_j\in [0, 1]
\]
Define $A^2= A\cdot A$, estimation of the ranks is simplified by the observation:
\[
\frac{1}{n}\mathbb{E}((A^2)_{ij}|Z) = \frac{1}{n}\sum_{k=1}^n z_i^T z_k z_k^T z_j = z_i^T \frac{\left(\sum_{k=1}^{n}z_{k}z_{k}^{T}\right)}{n}z_{j} = z_{i}^{T} \hat{\Sigma}_Z z_{j}
\]
Then finding the ranks of $(A^2)_{ij}$ we can find the ranks of $P_{ij}=z_i^T z_j$.
\end{frame}

% \begin{frame}{Rank estimation - weighted network}
% When does this work ? We need the matrix $\hat{\Sigma}$ to preserve the ordering, necessary conditions would be:
% \begin{itemize}
%     \item $\hat{\Sigma}$ is full rank. If one of the eigenvalues is zero one dimension would be "collapsed" and become irrelevant for the ordering.
% \end{itemize}

% \begin{itemize}
%     \item when the covariance matrix is diagonal such that $\hat{\Sigma}\to cI$ 
%     \item in simulation works also when $\Sigma$ is supposed to disto
% \end{itemize}
% \vspace{.5em}
% \end{frame}

% \begin{frame}{Ranking estimation - Binary network}
% For a RDPG model where $z_i\sim F$ such that $z_i^T z_j\in [0, 1]$: 
% \[
% A_{ij}\sim Bernoulli(z_i^T z_j)
% \]
% we have the same decomposition as before.
% \vspace{.5em}

% If we instead assume a inner product model with a link function $f$ we would have:
% \[
% \frac{1}{n}\mathbb{E}(A^2_{ij}) = \frac{1}{n}\sum_k f(x_i^T x_k) f(x_j^Tx_k)
% \]
% % This seems to work in practice. 

% % There may also be a connection with graphon distance as defined by \citet{zhang_estimating_2017}:
% % \[
% % d(i, j) = \max_{k\neq i,j} |\langle A_i-A_{j}, A_k\rangle|
% % \]
% \end{frame}

\section{Simulation results}

\begin{frame}{Simulations setup}
The model is made of 3 components:
\[
p(A, B) = \underbrace{p(A|Z) p(B|X)}_{\textbf{\color{red}likelihood}} \; \underbrace{p(Z) p(X)}_{\textbf{\color{green!50!black}marginals}} \; \underbrace{c(X, Z)}_{\textbf{\color{blue}copula}}
\]
The setup tested are:
\begin{itemize}
\item {\color{red}Likelihood}: Gaussian $A_{ij}\sim \mathcal{N}(z_i^T z_j, \sigma)$
\item {\color{green!50!black}Marginals}: $\mathcal{N}(0, 1)$, $t_5$ (df=5), $\chi^2_5$ (df=5), $Unif(-1, 1)$, Cauchy
\item {\color{blue}Copula}: Gaussian, Clayton, Gumbel, Student t (df=5), mixture of Gaussians
\end{itemize}
Results for archimedean and elliptical copulas are for $\rho=0.2$. Results for mixture of Gaussian are for $\rho=0.5$.
\end{frame}


\begin{frame}{Simulation Setup}
Models tested are:
\begin{itemize}
    \item Quadratic Assignment Procedure permutation test (QAP)
    \item Distance Correlation permutation test using MGC (DC)
    \item RV coefficient permutation test (RV)
    \item CvM permutation test (CVM)
\end{itemize}

\vspace{.5em}
For the methods requiring it, latent positions are estimated using:
\begin{itemize}
    \item \textit{Projected gradient descent} \citep{ma_universal_2020} for binary networks
    \item \textit{Adjacency Spectral Embedding} (ASE) for weighted networks
\end{itemize}
These algorithms are run using the "true" dimension $d$ (i.e. the $d$ used to generate the data).
\vspace{.5em}

% Note that in their work on Distance Correlation \citet{lee_network_2019} use diffusion maps to estimate latent positions. Here however we use the same estimator for every model to isolate the effect of the dependence metric choice.
\end{frame}

\begin{frame}{Type I error - Weighted network}
\vspace{-.5em}

\begin{figure}
    \centering
    \includegraphics[width=.6\linewidth]{figures/false_rejection_rate_weighted_network.png}
    \label{fig:placeholder}
\end{figure}
\vspace{-3em}
\end{frame}

% \begin{frame}{Power analysis}
% \begin{figure}
%     \centering
%     \includegraphics[width=\linewidth]{figures/blank_part1.png}
%     \label{fig:placeholder}
% \end{figure}
% \vspace{-2em}

% \begin{figure}
%     \centering
%     \includegraphics[width=0.9\linewidth]{figures/blank_part2.png}
%     \label{fig:placeholder}
% \end{figure}
% \end{frame}

% \begin{frame}{Power analysis - Weighted network}
% Results for $n=300, d=3, \sigma=3$.
% \begin{figure}
%     \centering
%     \includegraphics[width=\linewidth]{figures/first_piece.png}
%     \label{fig:placeholder}
% \end{figure}
% \vspace{-2em}

% \end{frame}

\begin{frame}{Power analysis - Weighted network}
Results for $n=300, d=3, \sigma=3$.
\begin{figure}
    \centering
    \includegraphics[width=\linewidth]{figures/second_piece.png}
    \label{fig:placeholder}
\end{figure}
\vspace{-2em}

\end{frame}

\begin{frame}{Power analysis - Weighted network}
Results for $n=300, d=3, \sigma=3$.
\begin{figure}
    \centering
    \includegraphics[width=\linewidth]{figures/third_piece.png}
    \label{fig:placeholder}
\end{figure}
\vspace{-2em}

\end{frame}

\begin{frame}{Power analysis - Weighted network}
Results for $n=300, d=3, \sigma=3$.
\begin{figure}
    \centering
    \includegraphics[width=\linewidth]{figures/fourth_piece.png}
    \label{fig:placeholder}
\end{figure}
\vspace{-2em}

\end{frame}




\begin{frame}
\vfill
\centering
{\Huge Thank you :)}

\vspace{0.5cm}
{\Large Questions?}
\vfill
\end{frame}

\begin{frame}[allowframebreaks]{Bibliography}    
\bibliography{references}
\end{frame}

\section{Appendix}
\subsection{More non-linear dependence}
\begin{frame}{More non-linear correlations from \citet{shen_distance_2020}}
To better visualise the issues in detecting non-linear dependency we use simulations settings from \citet{shen_distance_2020} who run comprehensive experiments aggregating 20 dependences from the literature (shown in the next slides). 
\vspace{.5em}

We then generate $X, Z$ and $A_{ij}\sim \mathcal{N}(z_i^T z_j, 1), B_{ij}\sim \mathcal{N}(x_i^T x_j, 1)$ for $z_i, x_i\in \mathbb{R}$. Latent positions era estimated using ASE.
% \begin{itemize}
%     \item \textbf{Weighted network}: $A_{ij}\sim \mathcal{N}(x_i^T x_j, 1)$ for $x_i\in \mathbb{R}$
%     \item \textbf{RDPG}: following \citet{lee_network_2019} 
%     \begin{enumerate}
%         \item generate $x_i\in \mathbb{R}$
%         \item normalise as $\tilde{x}_i= \frac{x_i - \min(x_i)}{\max x_i - \min x_i}$
%         \item draw $A_{ij}\sim Bern(x_i^T x_j)$
%     \end{enumerate}
% \end{itemize}
\end{frame}

\begin{frame}{Weighted network}
First 4 are monotonic relationships, here QAP and RV coefficient should work best.
\begin{figure}
\centering
\includegraphics[width=0.75\linewidth]{figures/visualise_latent_rdpg_part1.png}
\end{figure}
\vspace{-1.5em}

\begin{figure}
\centering
\includegraphics[width=\linewidth]{figures/power_weighted_latent_sim_first_batch_highlighted.png}
\end{figure}
\end{frame}

\begin{frame}[noframenumbering]{Weighted network}
First 4 are monotonic relationships, here QAP and RVcoefficient should work best.
\begin{figure}
\centering
\includegraphics[width=0.75\linewidth]{figures/visualise_latent_rdpg_part1.png}
\end{figure}
\vspace{-1.5em}

\begin{figure}
\centering
\includegraphics[width=\linewidth]{figures/power_weighted_latent_sim_first_batch.png}
\end{figure}
\end{frame}


\begin{frame}{Weighted network}
For small $n$ we see some methods working better than others 
\begin{figure}
\centering
\includegraphics[width=0.75\linewidth]{figures/visualise_latent_rdpg_part2.png}
\end{figure}
\vspace{-1.5em}

\begin{figure}
\centering
\includegraphics[width=\linewidth]{figures/power_weighted_latent_sim_second_batch.png}
\end{figure}
\end{frame}

\begin{frame}{Weighted network}
At $n=300$ we have clear separation between method that work and don't work for different scenarios.
\begin{figure}
\centering
\includegraphics[width=\linewidth]{figures/power_weighted_latent_sim_first_batch_n300.png}
\end{figure}
\vspace{-1.5em}

\begin{figure}
\centering
\includegraphics[width=\linewidth]{figures/power_weighted_latent_sim_second_batch_n300.png}
\end{figure}
\end{frame}

\begin{frame}{Future directions}
Future directions:
\begin{itemize}
\item formalise when CVM works and when it does not
\item formalise the binary case
\item possibly extend to more than two graphs to test a sort of "global null" 
\end{itemize}
\end{frame}

\subsection{Binary network simulations}
\begin{frame}{Type I error}

\begin{figure}
    \centering
    \includegraphics[width=1.05\linewidth]{figures/false_rejection_rate_binary_network.png}
    \label{fig:placeholder}
\end{figure}
    
\end{frame}

\begin{frame}{Power analysis - Binary network}
Results for $A_{ij}\sim Bern(f(x_i^T x_j)$ ($f(\cdot)$ logistic function) with $n=300, d=3$
\begin{figure}
    \centering
    \includegraphics[width=\linewidth]{figures/power_binary_part_1.png}
    \label{fig:placeholder}
\end{figure}
\vspace{-2em}

\begin{figure}
    \centering
    \includegraphics[width=\linewidth]{figures/power_binary_part_2.png}
    \label{fig:placeholder}
\end{figure}
\end{frame}

\subsection{RDPG simulations}

\begin{frame}{Random Dot Product Graph}
First 4 are monotonic relationships, here QAP and RVcoefficient should work best.
\begin{figure}
\centering
\includegraphics[width=0.75\linewidth]{figures/visualise_latent_rdpg_part1.png}
\end{figure}
\vspace{-1.5em}

\begin{figure}
\centering
\includegraphics[width=\linewidth]{figures/power_binary_latent_sim_first_batch.png}
\end{figure}
\end{frame}

\begin{frame}{Random Dot Product Graph}
\begin{figure}
\centering
\includegraphics[width=0.75\linewidth]{figures/visualise_latent_rdpg_part2.png}
\end{figure}
\vspace{-1.5em}

\begin{figure}
\centering
\includegraphics[width=\linewidth]{figures/power_binary_latent_sim_second_batch.png}
\end{figure}
\end{frame}

\begin{frame}{Random Dot Product Graph}
At $n=300$ we have clear separation between method that work and don't work for different scenarios.
\begin{figure}
\centering
\includegraphics[width=\linewidth]{figures/power_binary_latent_sim_first_batch_n300.png}
\end{figure}
\vspace{-1.5em}

\begin{figure}
\centering
\includegraphics[width=\linewidth]{figures/power_binary_latent_sim_second_batch_n300.png}
\end{figure}
\end{frame}


\begin{frame}{Why CVM does not work ?}
For both weighted and RDPG models we are using:
\[
\frac{1}{n}\mathbb{E}(A^{2}_{ij}|Z) = z_{i}^{T} \hat{\Sigma}_Z z_{j}
\]
The difference is that for RDPG we need to normalise latent positions to be in $[0, 1]$. This considerably reduces the tails making the ordering harder to recover. 
\end{frame}

\begin{frame}{Distortion - weighted case}

\begin{figure}
\centering
\includegraphics[width=\linewidth]{figures/visualise_latent_rdpg_gaussian.png}
\end{figure}
    
\end{frame}

\begin{frame}{Distortion - RDPG case}

\begin{figure}
\centering
\includegraphics[width=\linewidth]{figures/visualise_latent_rdpg_bernoulli.png}
\end{figure}
    
\end{frame}
\subsection{Differnt condition number}
\begin{frame}{Different condition number}
Weighted network, $n=300, d=3, \sigma=1$. Rows show different condition number of $\Sigma_{X}$.
\begin{figure}
\centering
\includegraphics[width=.9\linewidth]{figures/power_diff_variance.png}
\end{figure}
    
\end{frame}

\subsection{True vs estimated latent}
\begin{frame}{True vs estimated latent}
Weighted network, $n=100, d=3, \sigma=1$. 
\begin{figure}
\centering
\includegraphics[width=\linewidth]{figures/power_true_latent_1.png}
\end{figure}
\vspace{-1.5em}

\begin{figure}
\centering
\includegraphics[width=\linewidth]{figures/power_true_latent_2.png}
\end{figure}
    
\end{frame}

\end{document}